\begin{definition}[$k$-tensor]
  Um tensor de grau $k$ ou um  \textbf{k-tensor} é uma função real $k$-linear
  \begin{align*}
    T:\underbrace{V\times\cdots\times K}_{k-vezes}:\to \mathbb{R}.
  \end{align*}
  O espaço vetoril dos vetores $k$-tensores em $V$ é denotado por $T^{k}(V)$. 
\end{definition}
\begin{example}
  \phantom{X}
  \begin{itemize}
    \item $k=0$ temos que $T^{0}(V)=\mathbb{R}$.
    \item $k=1$ temos que $T^{1}(V)=V^{\ast}=\{T:V\to\mathbb{R}: T\text{ é linear}\}$.
    \item $k=2$ temos que $T^{2}(V)$ são os tensores bilineares como por exemploo o produto interno.
    \item $k=n$ temos que $\det\in T^{n}(V)$. 
  \end{itemize}
\end{example}
\begin{definition}[$k$-formas]
  As \textbf{$k$-formas} ou tensores alternados são tensores tais que
  \begin{align*}
    w(v_1,\cdots,v_i,\cdots,v_j,\cdots,v_n)=-w(v_1,\cdots,v_j,\cdots,v_i,\cdots,v_n).
  \end{align*}
  O subespaço vetorial das $k$-formas de um espaço vetorial $V$ será denotado por $\Lambda^{k}(V)$. 
\end{definition}
\begin{example}
  \phantom{X}
  \begin{itemize}
    \item $\Lambda^{0}(V)=T^{0}(V)=\mathbb{R}$.
    \item $\Lambda^{1}(V)=T^{1}(V)=V^{\ast}$.
    \item $\Lambda^{2}(V)\subseteq T^{2}(V)$.
  \end{itemize}
\end{example}
\begin{proposition}
  Seja $\sigma\in S_n$ uma permutação. Se $w\in\Lambda^{k}(V)$, então
  \begin{align*}
    w(v_{\sigma(1)},\cdots,v_{\sigma(k)})=(\sgn(\sigma))w(v_1,\cdots,v_k).
  \end{align*}
\end{proposition}
\begin{proposition}
  As seguintes afirmações são equivalentes
  \begin{enumerate}[i)]
    \item $T$ é um tensor alternado.
    \item $T(v_1,\cdots,v,\cdots,v,\cdots,v_k)=0$ para todo $v$.
    \item Se $v_i=v_{i+1}=v$ para algúm $i$, então $w(v_1,\cdots,v,v,\cdots,v_k)=0$.
    \item $T(v_1,\cdots,v_k)=0$ sempre que $v_1,\cdots,v_k$ são linearmente dependentes. 
  \end{enumerate}
\end{proposition}
\begin{proof} 
  $i)\to ii)$\\
  Sabemos que
  \begin{align*}
    T(v,v)&=-T(v,v)\\
    2T(v,v)&=0\\
    T(v,v)&=0.
  \end{align*}
  $ii)\to i)$\\
  Temos que pela linearidad de $T$ e a hipótese se cumpre que 
  \begin{align*}
    0=T(v_i+v_j,v_i+v_j)=T(v_i,v_i)+T(v_j,v_i)+T(v_i,v_j)+T(v_j,v_j).
  \end{align*}
  Mas como $T(v,v)=0$, então temos que $0=T(v_i,v_j)+T(v_j,v_i)$ o que implica que
  \begin{align*}
    T(v_i,v_j)=-T(v_j,v_i).
  \end{align*}
  $ii)\to iii)$\\
  Inmediato.\\
  $iii)\to ii)$\\
  Raciocinando por indução sobre $n$ em $v_i=v_{i+n}$ temos que pela hipótese o caso $n=1$ se cumpre.\\
  Agora, suponha se $v_i=v_{i+j}$ para algúm $i$ e para qualquer $j\in\{1,\cdots,r\}$, então vale que $T(v_1,\cdots,v_k)=0$. Vamos probar que isto implica que se $v_i=v_{i+r+1}$, então $T(v_1,\cdots,v_k)=0$.\\
  De fato temos que usando as hipóteses temos que como $v_i=v_{i+r+1}$, então
  \begin{align*}
    T(v_i,v_{i+r},v_{i+r+1})&=T(v_i,v_{i+r},v_{i+r+1}+v_{i+r})-\underbrace{T(v_i,v_{i+r},v_{i+r})}_{\to 0}\\
    &=\underbrace{T(v_i,v_{i+r}+v_{i+r+1},v_{i+r+1}+v_{i+r})}_{\to 0}-T(\underbrace{v_i}_{i-ésima},\underbrace{v_{i+r+1}}_{i+r-ésima},v_{i+r+1}+v_{i+r})\\
    &=0.
  \end{align*}
  $iv)\to ii)$\\
  Inmediato.\\
  $ii)\to iv)$\\
  Suponha que
  \begin{align*}
    v_i=\sum_{\substack{j=1\\j\neq i}}^{k}x^{j}v_j.
  \end{align*}
  Logo temos que pela linearidad de $T$ na $i$-ésima coordenada 
  \begin{align*}
    T(v_1,\cdots,v_k)=\sum_{\substack{j=1\\j\neq i}}^{k}x^{j}T(v_1,\cdots,v_j,\cdots,v_k)=0.
  \end{align*}
\end{proof}
\begin{theorem}[Existencia e unicidade do tensor determinante]
  Seja $V$ um espaço vetorial e $B=\{e_1,\cdots,e_n\}$ uma base de $V$.\\
  Existe uma única $n$-forma $\det$ tal que $\det(e_1,\cdots,e_n)=1$.\\
  Consequentemente temos que
  \begin{align*}
    \det(v_1,\cdots,v_n)=\sum_{\sigma\in S_n}(\sgn(\sigma))v_1^{\sigma(1)}\cdots v_{n}^{\sigma(n)}.
  \end{align*}
\end{theorem}
\begin{proof}
  Definamos $\det$ como
  \begin{align*}
    \det(v_1,\cdots,v_n)=\sum_{\sigma\in S_n}(\sgn(\sigma))v_1^{\sigma(1)}\cdots v_{n}^{\sigma(n)}.
  \end{align*}
  \textbf{Afirmação:}\\
  Temos que $\det$ é uma $n$-forma, ou seja é $n$-multilinear e alternada e $\det(e_1,\cdots,e_n)=1$.
  \begin{itemize}
    \item Note que $\det$ é multilinear, ja que o produto dos $v_i^{\sigma(i)}$ é multilinear.
    \item Note que se pegamos $\gamma=(i,j)$, então
      \begin{align*}
        \det(e_{\gamma(1)},\cdots,e_{\gamma_{\gamma(n)}})&=\sum_{\sigma\in S_n}(\sgn(\sigma))\delta_{1}^{\sigma(\gamma(1))}\cdots\delta_{n}^{\sigma(\gamma(n))}\\
        &=\sum_{\sigma\in S_n}(\sgn(\sigma))(\sgn(\gamma))\delta_{1}^{\sigma(1)}\cdots\delta_{n}^{\sigma(n)}\\
        &=-\det(e_1,\cdots,e_n).
      \end{align*}
    \item Note que
      \begin{align*}
        \det(e_1,\cdots,e_n)=\sum_{\sigma\in S_n}(\sgn(\sigma))(\sgn(\gamma))\delta_{1}^{\sigma(1)}\cdots\delta_{n}^{\sigma(n)}=1.
      \end{align*}
  \end{itemize}
  Pelo que podemos afirmar existencia.\\
  Para mostrar a unicidade temos que como $\det(e_1,\cdots,e_n)=1$, então dados vetores arbitrarios
  \begin{align*}
    v_j=\sum_{i=1}^{n}v_j^{i}e_i,\quad\text{ com }j\in\{1,\cdots,n\}.
  \end{align*}
  Aplicamos a multilinearidade de $\det$ e temos que
  \begin{align*}
    \det(v_1,\cdots,v_n)=\det\left( \sum_{i_1=1}^{n}v_{1}^{i_1}e_{i_1},\cdots, \sum_{i_n=1}^{n}v_{n}^{i_n}e_{i_n}\right)=\sum_{i_1=1}^{n}\cdots\sum_{i_n=1}^{n}v_{1}^{i_1}\cdots v_{n}^{i_n}\det(e_{i_1},\cdots,e_{i_n}).
  \end{align*}
  Agora, suponha que existem $p,q$ taes que $i_p=i_q$ com $p\neq q$, então como $\det$ é uma forma alternada, temos que $\det(e_{i_1},\cdots,e_{i_n})=0$, logo temos que
  \begin{align*}
    \det(v_1,\cdots,v_n)=\sum_{\sigma\in S_n}(\sgn(\sigma))v_1^{\sigma(1)}\cdots v_{n}^{\sigma(n)}.
  \end{align*}
\end{proof}
\begin{proposition}
  Se $\dim(V)=n$, então $\Lambda^{k}(v)=0$ para todo $k>n$. 
\end{proposition}
\begin{proof} 
  Note que como $k>n$, entao temos más componentes que vetores da base, logo se pegamos $w\in\Lambda^{k}(v)$ se cumpre que
  \begin{align*}
    w(e_1,\cdots,e_k)=0.
  \end{align*}
  Pois tem que existir algúm $i\neq j$ tal que $e_i=e_j$. 
\end{proof}
